В совместном MBAR по всем температурам и энергиям с использованием каждого двенадцатого кадра разность равна \(+17{,}11\)~К [16,08; 18,03]. Объединение только соседних с вершинами изотерм даёт 1262,33 и 1244,34~К, разность \(+17{,}99\)~К [16,21; 19,37]. При использовании каждого четвёртого кадра получены 1261,74 и 1244,20~К, разность \(+17{,}55\)~К. Для этой проверки блочный интервал повторно не оценивался. Разброс центральных оценок по способам обработки дополняет приведённые для каждого способа условные 95\%-е интервалы.

В общем поле \(\Delta\mu=3{,}545\)~эВ при 1243~К изменение разности полубольших потенциалов ограниченных \(Q\)-бассейнов \(\omega_{\alpha_2}-\omega_\alpha\) равно \(-0{,}388\)~мэВ/атом [\(-0{,}411\); \(-0{,}366\)]. Энергетически малая поправка заметно сдвигает вершину. Для \(w=0{,}35\) глубина минимума от слабой метастабильной моды здесь составляет всего \(0{,}32\,k_{\rm B}T\), для \(w=0{,}40\) --- \(2{,}07\,k_{\rm B}T\). Значения теплоты в таблице получены в вершине каждого ансамбля.
